\chapter{واجهات نظام التشغيل}
\label{chap:unix}

تتمثل الوظيفة الجوهرية لنظام التشغيل في مشاركة العتاد الفيزيائي للحاسوب بين عدة برامج،
وتوفير طبقة من الخدمات الأساسية تزيد من كفاءة وفائدة ما يقدمه العتاد الصلب بمفرده.
يتولى نظام التشغيل إدارة الموارد المادية منخفضة المستوى وبناء طبقة تجريد (Abstraction) تفصل البرمجيات عنها؛
فلا يحتاج تطبيق معالج النصوص—على سبيل المثال—إلى معرفة نوع قرص التخزين الصلب أو تفاصيل التعامل الكهربائي معه.
كما يتولى نظام التشغيل توزيع الموارد المادية بين برامج متعددة،
بحيث تُنفذ تلك البرامج بالتوازي في وقت واحد (أو تُعطي للمستخدم انطباعاً بالتنفيذ المتزامن).
وأخيراً، يتيح نظام التشغيل آليات محكمة ومُنظّمة لربط البرامج ببعضها،
مما يُمكنها من تبادل البيانات والتعاون المتبادل لإنجاز المهام.

يقدم نظام التشغيل خدماته للتطبيقات عبر واجهة برمجة منسقة (Interface).
ويرتكز تصميم واجهات نظم التشغيل الناجحة على معادلة هندسية دقيقة:
فمن جهة، يُستحسن أن تتسم الواجهة بالبساطة وصغر النطاق،
تسهيلاً للتحقق من صحة واثوقية التنفيذ البرمجي داخل النواة.
ومن جهة أخرى، قد يتجه المصمم إلى إغراق الواجهة بميزات معقدة لتلبية كافة متطلبات التطبيقات.
وتتجلى براعة التصميم في بناء واجهات تعتمد على عدد محدود من الآليات البسيطة،
ولكنها تتميز بالمرونة والقابلية للتركيب، مما يوفر قدرة عالية على العمومية والتطوير.

يقدم هذا الكتاب نظام التشغيل \texttt{xv6} كنماذج تطبيقي وملموس لتوضيح المفاهيم الأساسية لنظم التشغيل.
يقدم نظام \texttt{xv6} الواجهات والخدمات الجوهرية التي جاء بها نظام Unix التاريخي،
الذي ابتكره كين طومسون ودنيس ريتشي~\cite{unix}، كما يحاكي الهيكلية الداخلية لنواة Unix.
يوفر Unix واجهة محكمة تتصف آلياتها بالتكامل وقابلية التركيب،
مما يمنح النظام مرونة وعمومية مدهشتين.
وقد حققت هذه الواجهة نجاحاً واسعاً حتى أصبحت النواة الأساسية للنظم الحديثة—مثل
BSD و Linux و macOS و Solaris و Microsoft Windows إلى حد ما.
ولذلك، فإن دراسة نظام \texttt{xv6} واستيعاب آلياته تشكل مدخلاً رئيساً لفهم أنظمة التشغيل الحديثة بشتى أنواعها.

\begin{figure}[t]
\centering
\includegraphics[width=0.8\textwidth]{fig/os}
\caption{نواة نظام التشغيل مع عمليتين مستخدمين.}
\label{fig:os}
\end{figure}

كما يتضح من الشكل~\ref{fig:os}، يتبنى \texttt{xv6} الهيكلية التقليدية القائمة على وجود \indextext{النواة (Kernel)}،
وهي برنامج متميز يعمل بوضع المشرف ويملك صلاحيات كاملة على العتاد لتوفير الخدمات للعمليات الجارية.
تُسمى كل وحدة تنفيذ جارية \indextext{عملية (Process)}،
وتخصص النواة لكل عملية مساحة ذاكرة مستقلة تضم تعليمات البرنامج، وبياناته، ومكدس الاستدعاءات (Stack).
تحتوي التعليمات على الشفرة التنفيذية، وتُمثل البيانات المتغيرات والمساحات التخزينية،
بينما ينظم المكدس استدعاءات الدوال والتعامل مع المتغيرات المحلية.

تضم الحاسوب المعتاد عدة عمليات جارية بالتوازي، في حين توجد نواة واحدة فقط تتولى الإدارة.
وعندما تحتاج العملية إلى الحصول على إحدى خدمات النواة، فإنها تصدر \indextext{استدعاء نظام (System Call)}،
وهو واجهة برمجية محددة تدعمها النواة.
ينتقل التحكم عند صدور استدعاء النظام إلى النواة، فتنفذ النواة الخدمة المطلوبة ثم تعود بالنتيجة إلى العملية.
وبذلك يتناوب التنفيذ بين \indextext{مساحة المستخدم (User Space)} و\indextext{مساحة النواة (Kernel Space)}.

يطلق على مجموعة استدعاءات النظام التي تتيحها النواة للتطبيقات اسم الواجهة الموجهة للتطبيقات.
ويوضح الجدول~\ref{t:sys} كافة استدعاءات النظام التي يوفرها نظام \texttt{xv6}.

\begin{table}[t]
\centering
\begin{tabular}{ll}
\toprule
\textbf{استدعاء النظام} & \textbf{الوصف والوظيفة} \\
\midrule
\texttt{fork()} & إنشاء عملية جديدة بذاكرة مطابقة للعملية الأب. \\
\texttt{exit(status)} & إنهاء العملية الحالية وتمرير حالة الخروج للعملية الأب. \\
\texttt{wait(\&status)} & الانتظار حتى تنتهي إحدى العمليات الأبناء واسترجاع حالة خروجها. \\
\texttt{pipe(p)} & إنشاء أنبوب لنقل البيانات وإعادة واصفات ملفات القراءة والكتابة. \\
\texttt{write(fd, buf, n)} & كتابة \texttt{n} بايتاً من المؤشر \texttt{buf} إلى واصف الملف \texttt{fd}. \\
\texttt{read(fd, buf, n)} & قراءة \texttt{n} بايتاً إلى المؤشر \texttt{buf} من واصف الملف \texttt{fd}. \\
\texttt{close(fd)} & إغلاق واصف الملف المحدد وتفريغ القناة. \\
\texttt{kill(pid)} & إنهاء العملية المحددة بمعرف \texttt{pid}. \\
\texttt{exec(filename, *argv)} & تحميل ملف تنفيذي وتنفيذه بالوسائط الممررة. \\
\texttt{open(filename, flags)} & فتح ملف وتحديد نمط الوصول وإعادة واصف ملف. \\
\texttt{mknod(dev, main, minor)} & إنشاء ملف جهاز خاص في نظام الملفات. \\
\texttt{unlink(filename)} & إزالة ملف أو رابط صلب من نظام الملفات. \\
\texttt{fstat(fd, status)} & استرجاع المعلومات والبيانات الوصفية عن الملف. \\
\texttt{link(path1, path2)} & إنشاء رابط صلب جديد يُشير إلى ملف موجود. \\
\texttt{mkdir(dirname)} & إنشاء مجلد جديد بداخل نظام الملفات. \\
\texttt{chdir(dirname)} & تغيير المجلد الحالي للعملية المداعية. \\
\texttt{dup(fd)} & تكرار واصف ملف قائم وتعيين أحدث واصف متاح له. \\
\bottomrule
\end{tabular}
\caption{استدعاءات النظام الأساسية في نظام التشغيل \texttt{xv6}.}
\label{t:sys}
\end{table}

تُنفذ النواة خدماتها باستغلال آليات العتاد الصلب لحماية المساحات التخزينية وتحديد صلاحيات التنفيذ.
يُنفذ برنامج المستخدم بوضع غير مميز (User Mode)، يمنعه العتاد فيه من الوصول المباشر إلى الذاكرة المخصصة للأجهزة أو النواة.
بينما تُنفذ النواة بوضع المشرف (Supervisor Mode)، مما يُتيح لها التعامل المباشر مع جميع الأجهزة وسجلات العتاد.

يُعد \indextext{مغلف الأوامر (Shell)} برنامج مستخدم كلاسيكي يتولى قراءة الأوامر من المستخدم وتنفيذها.
ولا يُعتبر الشل جزءاً داخل النواة، بل هو تطبيق عادي يبرز مرونة واجهة استدعاءات النظام وقوتها.

\section{العمليات والذاكرة}

تُعتبر \indextext{العملية (Process)} في \texttt{xv6} وحدة العزل والتنفيذ الأساسية؛
حيث يمنع نظام العزل عمليةً ما من تخريب ذاكرة عملية أخرى أو الوصول إليها دون إذن.
تتكون مساحة عنوان العملية من الذاكرة التي تحوي شفرة البرنامج، والبيانات، ومكدس التنفيذ.
يستعرض هذا الفصل كيف ينشئ \texttt{xv6} العمليات ويدير ذاكرتها عبر استدعاءات النظام المخصصة.

تتيح استدعاءات \texttt{fork} و \texttt{exit} و \texttt{wait} إدارة دورة حياة العمليات.
أنشئ العملية الأب عملية ابن جديدة عبر استدعاء \texttt{fork()}،
فتقوم النواة بنسخ كامل مساحة العنوان (الشفرة والبيانات والمكدس) للعملية الأصلية إلى مساحة عملية جديدة.
يرجع استدعاء \texttt{fork} في كلتا العمليتين:
فيرجع بمعرّف العملية الابن (\texttt{PID}) للعملية الأب، بينما يرجع بالقيمة \texttt{0} داخل العملية الابن نفسها.

يوضح مقطع البرنامج التالي المكتوب بلغة C الاستخدام المباشر لـ \texttt{fork}:

\begin{lstlisting}
int pid = fork();
if(pid > 0){
  printf("parent: child=%d\n", pid);
  pid = wait((int *) 0);
  printf("child %d is done\n", pid);
} else if(pid == 0){
  printf("child: exiting\n");
  exit(0);
} else {
  printf("fork error\n");
}
\end{lstlisting}

ينهي استدعاء \texttt{exit(status)} العملية المداعية ويُحرر مواردها، وتتسلم العملية الأب حالة الخروج عند استدعاء \texttt{wait(\&status)}.
يتوقف الأب عند تعليق التنفيذ في \texttt{wait} حتى تنتهي إحدى العمليات الأبناء، مما يحقق التزامن بين الأب وابنه.

أما استدعاء \texttt{exec(filename, argv)} فيستبدل الذاكرة الحالية للعملية المداعية ببرنامج جديد يُقرأ من الملف التنفيذي المحدد.
يقرأ \texttt{exec} الشفرة والبيانات ويُجهز مكدساً جديداً يحتوي على الوسائط \texttt{argv}، ثم يبدأ التنفيذ من نقطة الدخول الرئسية للبرنامج الجديد.
إذا نجح \texttt{exec}، فلا يعود إلى الشفرة القديمة، لأن مساحة الذاكرة القديمة تكون قد استبدلت تماماً.

يقوم الشل الأساسي في \texttt{xv6} بإنشاء العمليات وتنفيذ الأوامر عبر حلقة تكرارية:
يقرأ الشل السطر المدخل من المستخدم، ثم يستدعي \texttt{fork} لإنشاء عملية ابن،
ثم تستدعي العملية الابن \texttt{exec} لتنفيذ الأمر المطلوب، بينما يستدعي الأب \texttt{wait} بانتظار اكتمال تنفيذ الأمر.

\section{الإدخال والإخراج وواصفات الملفات}

يمثل \indextext{واصف الملف (File Descriptor)} عدداً صحيحاً صغيراً يمثل قناة مفتوحة بواسطة العملية لقراءة البيانات أو كتابتها.
تستطيع العملية قراءة البيانات أو كتابتها عبر الواصف بغض النظر عما إذا كان هذا الواصف متصلاً بملف على القرص، أو بذاكرة أنبوبية، أو بجهاز إدخال/إخراج كالشاشة واللوحة.
يُسمى هذا التجريد بـ "عمومية نظام الملفات"، حيث تتعامل كافة وسائل الإدخال والإخراج بنفس الواجهة الموحدة: \texttt{read} و \texttt{write}.

تخصص النواة بصفة افتراضية لكل عملية جديدة ثلاثة واصفات ملفات مفتوحة عند بدء التنفيذ:
\begin{itemize}
\item الواصف \texttt{0}: الإدخال القياسي (Standard Input)، لقراءة المدخلات من اللوحة.
\item الواصف \texttt{1}: الإخراج القياسي (Standard Output)، لكتابة مخرجات النص على الشاشة.
\item الواصف \texttt{2}: إخراج الأخطاء القياسي (Standard Error)، لطباعة رسائل الأخطاء على الشاشة.
\end{itemize}

يتيح استدعاء \texttt{read(fd, buf, n)} قراءة عدد من البايتات يصل إلى \texttt{n} من واصف الملف \texttt{fd} إلى المؤشر \texttt{buf}.
ويرجع الاستدعاء بعدد البايتات المقروءة فعلياً. إذا وصل المخصّص إلى نهاية الملف (EOF)، يرجع الاستدعاء بالقيمة \texttt{0}.
بينما يكتب الاستدعاء \texttt{write(fd, buf, n)} عدد \texttt{n} بايتاً من \texttt{buf} إلى واصف الملف \texttt{fd}، ويرجع بعدد البايتات المكتوبة.

تعتمد النواة قاعدة صارمة ومهمة عند تخصيص واصفات الملفات الجديدة:
\textit{"تختار النواة دائماً أصغر رقم واصف ملف غير مستخدم في العملية"}.
يستغل الشل هذه القاعدة البسيطة لإنجاز عملية \indextext{إعادة توجيه الإدخال والإخراج (I/O Redirection)}.

لإعادة توجيه المخرجات إلى ملف عبر الأمر \texttt{cat < input.txt > output.txt}، يقوم الشل بالخطوات التالية بداخل العملية الابن قبل استدعاء \texttt{exec}:
1. يستدعي \texttt{close(1)} لإغلاق الواصف القياسي للإخراج.
2. يستدعي \texttt{open("output.txt", O_CREATE|O_WRONLY)} لفتح ملف جديد.
وبما أن الواصف \texttt{1} أصبح فارغاً، فتضمن قاعدة النواة أن يُخصص الرقم \texttt{1} للملف المغلق حديثاً.
وعندما تستدعي العملية الابن برنامج \texttt{cat} عبر \texttt{exec}، فإن أي كتابة على الواصف القياسي \texttt{1} ستتجه مباشرة إلى الملف \texttt{output.txt}.

يتيح استدعاء \texttt{dup(fd)} تكرار واصف ملف قائم، فتُخصص النواة أصغر واصف متاح يُشير إلى نفس مدخل الملف المفتوح.
ويتشارك الواصف الأصلي والواصف المكرر نفس مؤشر التمرير (File Offset)، مما يعني أن الكتابة عبر أحدهما تحدّث موضع الكتابة للآخر فوراً.

\section{الأنابيب}

يُمثل \indextext{الأنبوب (Pipe)} قناة بيانات ثنائية الأطراف في مساحة النواة، تُتيح الاتصال المباشر بين عمليتين.
يُنشأ الأنبوب عبر استدعاء النظام \texttt{pipe(p)}، حيث يستقبل مصفوفة من عنصرين \texttt{p[2]}:
فتخصص النواة الواصف \texttt{p[0]} لعملية القراءة، والواصف \texttt{p[1]} لعملية الكتابة.

يوضح المثال التالي تواصل عمليتين عبر الأنبوب:

\begin{lstlisting}
int p[2];
char *argv[2];
argv[0] = "wc";
argv[1] = 0;

pipe(p);
if(fork() == 0) {
  close(0);
  dup(p[0]);
  close(p[0]);
  close(p[1]);
  exec("/bin/wc", argv);
} else {
  close(p[0]);
  write(p[1], "hello world\n", 12);
  close(p[1]);
}
\end{lstlisting}

يحقق الأنبوب مزايا هامة بالمقارنة مع استخدام الملفات المؤقتة على القرص:
\begin{enumerate}
\item \textbf{التنظيف التلقائي}: يُفرغ النظام الأنبوب ويُحرر ذاكرته فور إغلاق جميع واصفاته.
\item \textbf{القدرة على نقل سعات لا نهائية}: يُمكن تدفق جيجابايت من البيانات عبر الأنبوب دون الحاجة لمساحات قرصية كبيرة.
\item \textbf{التزامن الضمني}: إذا كان الأنبوب فارغاً، يتوقف استدعاء \texttt{read} محظوراً بانتظار كتابة بيانات؛ وإذا امتص الأنبوب سعته القصوى، يتوقف \texttt{write} حتى تقرأ العملية الأخرى البيانات.
\end{enumerate}

\section{نظام الملفات}

يوفر نظام الملفات في \texttt{xv6} المجلدات الشجرية والملفات المنظمة لتخزين البيانات واسترجاعها بصفة دائمة.
تضم مساحة نظام الملفات نوعين أساسيين:
\begin{itemize}
\item \textbf{الملفات العادية (Regular Files)}: كتل بيانات غير مشكلة تُقرأ وتُكتب كمتتالية من البايتات.
\item \textbf{المجلدات (Directories)}: ملفات خاصة تحتوي على قائمة بأسماء الملفات المربوطة بأرقام العقد الفهرسية.
\end{itemize}

تُمثل \indextext{العقدة الفهرسية (Inode)} الهيكل الأساسي المسؤول عن تمثيل الملف على القرص؛
حيث تحمل البيانات الوصفية للملف (مثل نوع الملف، وحجمه بالبايتات، وعدد الروابط الصلبة المشارة إليه، وقائمة كتل القرص التي تتضمن البيانات).
لا يحتوي اسم الملف على بيانات الملف نفسها، بل يُشير الاسم بداخل المجلد إلى رقم العقدة الفهرسية للملف.

تتيح استدعاءات النظام التالية إدارة نظام الملفات:
\begin{itemize}
\item \texttt{mkdir("dirname")}: إنشاء مجلد جديد.
\item \texttt{mknod("devpath", major, minor)}: إنشاء ملف جهاز خاص للتعامل مع العتاد الفيزيائي.
\item \texttt{fstat(fd, \&st)}: استرجاع معلومات العقدة الفهرسية وحالتها بداخل هيكل البيانات \texttt{struct stat}.
\item \texttt{link("oldpath", "newpath")}: إنشاء رابط صلب جديد يُشير إلى نفس رقم العقدة الفهرسية للملف القديم.
\item \texttt{unlink("path")}: إزالة اسم الملف من المجلد وتقليل عدد الروابط بداخل العقدة الفهرسية.
\end{itemize}

لا تُمسح بيانات العقدة الفهرسية من القرص فور استدعاء \texttt{unlink} إلا إذا وصل عدد الروابط المشارة إليها إلى الصفر وأغلقت كافة العمليات واصفات الملفات المفتوحة لهذه العقدة.

\section{العالم الحقيقي}

تعتبر الواجهات والخدمات المقدمة في \texttt{xv6} تطبيقاً رصيداً ومحاكياً لمعيار \indextext{POSIX} التاريخي في أنظمة Unix.
تحتوي الأنظمة التجارية والحديثة مثل Linux و FreeBSD و macOS على العديد من استدعاءات النظام المتقدمة (مثل دعم الأنسجة المقيدة، والذاكرة المشاركة، والتحكم بالشبكات)،
ولكن يظل النمط الهيكلي القائم على التفرع عبر \texttt{fork} والتبديل عبر \texttt{exec} والتعامل مع الملفات عبر الواصفات هو الحجر الأساس في كافة الأنظمة المعتمدة على معيار Unix.

\section{تمارين}

\begin{enumerate}
\item اكتب برنامجاً صغيراً بداخل \texttt{xv6} يُسمى \texttt{pingpong} يستخدم الأنابيب لإرسال بايت واحد بين عمليتين (الأب والابن) بالتناوب وطباعة رسائل التأكيد.
\item اكتب نسخة مصغرة من أداة \texttt{find} تبحث بداخل المجلدات الشجرية عن ملف يحمل اسماً معيناً وتطبع مساره الكامل.
\item اكتب أداة \texttt{xargs} تستقبل السطور من الإدخال القياسي وتنفذ أمراً ممرراً عبر الوسائط بداخل عملية ابن لكل سطر.
\end{enumerate}
